\section{Problem formulation}
\label{sec:problem}

\subsection{The agent loop and its feedback record}

An agent solving a task $\tau$ interacts with an environment over discrete
steps. At step $t$ the harness assembles a context $\ctxpre_t$ from the system
prompt, the goal, and a transcript of the previous steps; the policy $\pi$
samples an action $a_t \sim \pi(\cdot \mid \ctxpre_t)$; the environment executes
it and returns an observation $o_t$ together with a success flag. Writing
$\oplus$ for appending a turn pair to the transcript,
\begin{equation}
\ctxpre_{t+1} \;=\; \ctxpre_t \oplus (a_t, o_t).
\label{eq:loop}
\end{equation}

Equation~\eqref{eq:loop} is the design decision this paper is about. Almost every
agent framework implements it literally: the action the model wrote is appended
\emph{verbatim}, and so is whatever the runtime printed in response. When the
action fails, the transcript therefore contains the failed action's exact token
sequence, immediately followed by text asserting that it did not work.

The intent is obvious and, for capable models, well supported: the error message
is information, and a model that reads it should avoid the action that produced
it. We take that intent seriously enough to measure it.

\subsection{Corrective gain}

Fix a decision point: a context $\ctxpre$ before an attempt, an action
$\afail$ that will fail there, its observation $o^{\times}$, and the correct
action $\agold$. Define the \emph{corrective gain}
\begin{equation}
\Gain(\afail) \;=\;
\log \pi(\afail \mid \ctxpre)
\;-\;
\log \pi\!\left(\afail \mid \ctxpre \oplus (\afail, o^{\times})\right),
\label{eq:gain}
\end{equation}
where $\log \pi(a \mid c) = \sum_{i} \log \pi(a_i \mid c, a_{<i})$ is the summed
token log-probability of the action string. Both terms score \emph{the same
string}, so their difference is exactly the change in the log-odds of writing
that action again, and it is invariant to how long the action is or how the
tokeniser splits it.

$\Gain > 0$ is the behaviour the harness is designed to produce: after the
failure is recorded, the failed action is less likely. $\Gain < 0$ means the
record made the model \emph{more} likely to repeat the action; we call this
\emph{feedback inversion}.

\subsection{Decomposing the gain}

Appending $(\afail, o^{\times})$ does two things at once. It places the token
sequence of $\afail$ into the context, where the model's copying machinery can
reach it \citep{olsson2022induction,xu2022breaktheloop}; and it asserts that this
action failed. These pull in opposite directions, and $\Gain$ only reports the
net result.

We separate them with counterfactual observations. Let $o^{\varnothing}$ be a
\emph{valence-free} observation, which acknowledges the call but reports neither
success nor failure, and let $o^{\checkmark}$ be a counterfactual observation
reporting that the same call succeeded. Then
\begin{align}
\copyterm &= \log \pi\!\left(\afail \mid \ctxpre \oplus (\afail, o^{\varnothing})\right) \notag \\
          &\quad - \log \pi(\afail \mid \ctxpre), \label{eq:copy}\\
\semterm  &= \log \pi\!\left(\afail \mid \ctxpre \oplus (\afail, o^{\times})\right) \notag \\
          &\quad - \log \pi\!\left(\afail \mid \ctxpre \oplus (\afail, o^{\varnothing})\right), \label{eq:sem}
\end{align}
and by construction
\begin{equation}
-\Gain(\afail) \;=\; \copyterm + \semterm .
\label{eq:decomp}
\end{equation}
$\copyterm$ isolates what merely \emph{having written the action} does, holding
the feedback content neutral; $\semterm$ isolates what \emph{marking it as
failed} adds, holding the surface form fixed. We additionally report the
polarity contrast
\begin{equation}
\begin{split}
\Delta_{\mathrm{pol}} = {}& \log \pi\!\left(\afail \mid \ctxpre \oplus (\afail, o^{\times})\right) \\
                        & - \log \pi\!\left(\afail \mid \ctxpre \oplus (\afail, o^{\checkmark})\right),
\end{split}
\label{eq:pol}
\end{equation}
which compares two contexts that are identical except for whether the
observation says the call worked. A model that does not read the error at all
has $\Delta_{\mathrm{pol}} = 0$.

\subsection{What a fix would have to look like}

Equation~\eqref{eq:decomp} makes the design space concrete. If $\semterm$ is the
problem, the model does not understand the error, and the remedy is a better
model or a clearer message. If $\copyterm$ is the problem, no amount of message
engineering helps, because the damage is done by the presence of the string
rather than by anything said about it; the remedies are then \emph{structural}:
remove the surface form from the context, or make it unreachable at decoding
time. Section~\ref{sec:agent} evaluates both, against the natural-language
alternative of simply instructing the model not to repeat itself.
